\chapter{2012:Robert J. Lefkowitz and Brian Kobilka}

\label{chap:chemistry2012}

The Nobel Prize in Chemistry for the year 2012 was awarded jointly to Robert J. Lefkowitz and Brian Kobilka.

\textbf{Motivation:}

for studies of G-protein-coupled receptors

\section{Robert J. Lefkowitz}

\subsection*{Early Life and Education}

Robert J. Lefkowitz was born in 1943 in New York City, New York, USA.

\subsection*{Career and Major Research}

Lefkowitz earned his medical degree from Columbia University in 1966. He became professor of medicine and biochemistry at Duke University and an investigator of the Howard Hughes Medical Institute. He identified the receptors for adrenaline as G-protein-coupled receptors, cloned several of them, and discovered the arrestin proteins that regulate them.

\subsection*{Nobel Prize-winning Work}

The work for which Robert J. Lefkowitz was awarded the Nobel Prize in Chemistry in 2012 was described by the Nobel Committee as: "for studies of G-protein-coupled receptors".

Lefkowitz showed that a large family of receptors acts through G proteins, explaining how hormones and neurotransmitters signal across the cell membrane, and he uncovered the molecules that switch these signals off.

\subsection*{Significance and Impact}

G-protein-coupled receptors mediate the responses of cells to light, odor, and most hormones, and they are the targets of a large fraction of approved drugs. Lefkowitz's work established how these receptors recognize their signals and are regulated.

\subsection*{Later Career and Honors}

Lefkowitz continues his research at the Duke University School of Medicine and the Howard Hughes Medical Institute.

\subsection*{Legacy}

His discovery of arrestins and his account of receptor regulation are central to pharmacology, and his laboratory trained a generation of researchers in receptor biology.

\subsection*{Modern Developments}

Lefkowitz's discovery of the beta-adrenergic receptor and his characterization of G protein-coupled receptors (GPCRs) revealed the largest family of drug targets in medicine. Over a third of all approved drugs act on GPCRs, including treatments for asthma, heart disease, and hypertension, and his work founded the modern field of receptor pharmacology.

\subsection*{Selected Publications}

"Beta-Arrestins: A New Family of Regulators of G-Protein-Coupled Receptors" (Journal of Biological Chemistry, 1990)

\clearpage

\section{Brian Kobilka}

\subsection*{Early Life and Education}

Brian Kobilka was born in 1955 in Little Falls, Minnesota, USA.

\subsection*{Career and Major Research}

Kobilka earned a bachelor's degree from the University of Minnesota Duluth in 1977 and a medical degree from Washington University in St. Louis in 1981. He worked in Robert Lefkowitz's laboratory at Duke University and then became a professor at Stanford University, where in 2007 he determined the crystal structure of the beta-2 adrenergic receptor.

\subsection*{Nobel Prize-winning Work}

The work for which Brian Kobilka was awarded the Nobel Prize in Chemistry in 2012 was described by the Nobel Committee as: "for studies of G-protein-coupled receptors".

Kobilka obtained the first high-resolution structure of a G-protein-coupled receptor bound to a diffusible signal, and later captured the receptor in action with the G protein it activates.

\subsection*{Significance and Impact}

The structures revealed at atomic detail how a drug molecule binds the receptor and how the receptor changes shape to transmit its signal, explaining how GPCRs work as molecular switches.

\subsection*{Later Career and Honors}

Kobilka continues his research on receptor structure at the Stanford University School of Medicine.

\subsection*{Legacy}

Kobilka's structural work provided the framework for understanding GPCR activation and has aided the design of drugs that act on this family of receptors.

\subsection*{Modern Developments}

Kobilka's determination of the three-dimensional structure of the beta2-adrenergic receptor in its active state revealed how GPCRs are activated and how drugs stabilize them. His structures underpin modern structure-based drug design for receptors that are targets of a large fraction of pharmaceuticals, including beta-blockers and antihistamines.

\subsection*{Selected Publications}

"Crystal Structure of the Human $\beta_{2}$ Adrenergic G-Protein-Coupled Receptor" (Nature, 2007)

\clearpage

\section*{Chapter Summary}

The 2012 prize recognized the studies of G-protein-coupled receptors by Lefkowitz and Kobilka, which explained how these receptors transmit signals into cells. Their work is the foundation of much of modern pharmacology, since GPCRs are targets of many medicines.
